% SPDX-FileCopyrightText: 2026 Will Cook
% SPDX-License-Identifier: CC-BY-4.0
\documentclass[11pt]{article}
\usepackage[T1]{fontenc}
\usepackage[utf8]{inputenc}
\usepackage[a4paper,margin=1in]{geometry}
\usepackage{amsmath,amssymb,amsthm,booktabs,array}
\usepackage[dvipsnames]{xcolor}
\usepackage[colorlinks=true,linkcolor=black,urlcolor=RoyalBlue,citecolor=black]{hyperref}
\hypersetup{
  pdftitle={Finite-Difference Certificates and a Phase-Separation Characterisation for Erdős Problem 249},
  pdfauthor={Will Cook},
  pdfsubject={Companion note on transport, curvature, and local obstructions; specialist companion to the Erdős 249/257 gateway paper},
}
\newcommand{\commit}{c24e2c2a781aa4a104d10884bacae52ef35cff30}
\newcommand{\repobase}{https://github.com/wcook04/plectis-lean-erdos249-257/blob/\commit}
\newcommand{\PK}{Erdos249257}
% SPDX-FileCopyrightText: 2026 Will Cook
% SPDX-License-Identifier: CC-BY-4.0
% Generated by scripts/build_paper_module_aliases.py; do not edit.
% Each CamelCase component contributes at most its first three characters.
\newcommand{\modulesigil}[1]{\csname modulesigil@#1\endcsname}
\expandafter\def\csname modulesigil@ActualForeignResidueProjection.lean\endcsname{ActForResPro}
\expandafter\def\csname modulesigil@AdelicHeightObstruction.lean\endcsname{AdeHeiObs}
\expandafter\def\csname modulesigil@BooleanMobiusCarry.lean\endcsname{BooMobCar}
\expandafter\def\csname modulesigil@CampbellShiftSynchronization.lean\endcsname{CamShiSyn}
\expandafter\def\csname modulesigil@CarrySurvivorExtinction.lean\endcsname{CarSurExt}
\expandafter\def\csname modulesigil@CertificateKernel.lean\endcsname{CerKer}
\expandafter\def\csname modulesigil@DiagonalFlexibleOddWindowAffine.lean\endcsname{DiaFleOddWinAff}
\expandafter\def\csname modulesigil@DiagonalFlexibleOddWindowSupply.lean\endcsname{DiaFleOddWinSup}
\expandafter\def\csname modulesigil@DiagonalFreshLossBridge.lean\endcsname{DiaFreLosBri}
\expandafter\def\csname modulesigil@DiagonalPincerCertificateT64Endpoint.lean\endcsname{DiaPinCerT64End}
\expandafter\def\csname modulesigil@DiagonalPincerCertificatesT64.lean\endcsname{DiaPinCerT64}
\expandafter\def\csname modulesigil@DiagonalPincerDecomposition.lean\endcsname{DiaPinDec}
\expandafter\def\csname modulesigil@DiagonalPincerPrimeCertificates/ClosureT64.lean\endcsname{DiaPinPriCerT64}
\expandafter\def\csname modulesigil@FreshPrimeDeficitDecomposition.lean\endcsname{FrePriDefDec}
\expandafter\def\csname modulesigil@FullTargetPrimeAdjunctionNoGo.lean\endcsname{FulTarPriAdjNoGo}
\expandafter\def\csname modulesigil@GapFareyBound.lean\endcsname{GapFarBou}
\expandafter\def\csname modulesigil@GcdMomentCalculus.lean\endcsname{GcdMomCal}
\expandafter\def\csname modulesigil@GenericTailOrbitRigidity.lean\endcsname{GenTaiOrbRig}
\expandafter\def\csname modulesigil@GreedyAchievementSet.lean\endcsname{GreAchSet}
\expandafter\def\csname modulesigil@LcmConeFlatness.lean\endcsname{LcmConFla}
\expandafter\def\csname modulesigil@LcmConeNonflat.lean\endcsname{LcmConNon}
\expandafter\def\csname modulesigil@LcmDiagonalReduction.lean\endcsname{LcmDiaRed}
\expandafter\def\csname modulesigil@MersenneLambertLadder.lean\endcsname{MerLamLad}
\expandafter\def\csname modulesigil@MersenneShadowDenominatorGrowth.lean\endcsname{MerShaDenGro}
\expandafter\def\csname modulesigil@MersenneTailAtoms.lean\endcsname{MerTaiAto}
\expandafter\def\csname modulesigil@MobiusSignSupportNoGo.lean\endcsname{MobSigSupNoGo}
\expandafter\def\csname modulesigil@PowerTwoBitLift.lean\endcsname{PowTwoBitLif}
\expandafter\def\csname modulesigil@PowerTwoCenteredBitLift.lean\endcsname{PowTwoCenBitLif}
\expandafter\def\csname modulesigil@PrimitiveDeterminantLift.lean\endcsname{PriDetLif}
\expandafter\def\csname modulesigil@PrimitiveWeightCertificate.lean\endcsname{PriWeiCer}
\expandafter\def\csname modulesigil@RationalSupportCarrySkeleton.lean\endcsname{RatSupCarSke}
\expandafter\def\csname modulesigil@ResidualGaugeObstruction.lean\endcsname{ResGauObs}
\expandafter\def\csname modulesigil@SignedQMomentObstruction.lean\endcsname{SigQMomObs}
\expandafter\def\csname modulesigil@SquareCRTCube.lean\endcsname{SquCRTCub}
\expandafter\def\csname modulesigil@SquaredMersenneDiagonalEnclosure.lean\endcsname{SquMerDiaEnc}
\expandafter\def\csname modulesigil@SternBrocotRunGeometry.lean\endcsname{SteBroRunGeo}
\expandafter\def\csname modulesigil@SublogDivisorCoverage.lean\endcsname{SubDivCov}
\expandafter\def\csname modulesigil@TotientCarryKernelRigidity.lean\endcsname{TotCarKerRig}
\expandafter\def\csname modulesigil@TotientMahlerDefect.lean\endcsname{TotMahDef}
\expandafter\def\csname modulesigil@TotientTailPeriodKiller.lean\endcsname{TotTaiPerKil}

\newcommand{\lref}[3]{\href{\repobase/\PK/#1\#L#2}{Lean proof}}
\newcommand{\lword}[4]{\href{\repobase/\PK/#1\#L#2}{#4}}
\emergencystretch=2em
\newtheorem{theorem}{Theorem}[section]
\newtheorem{proposition}[theorem]{Proposition}
\theoremstyle{remark}\newtheorem*{remark}{Remark}
\newcommand{\Z}{\mathbb Z}
\newcommand{\ph}{\varphi}
\title{Finite-Difference Certificates and a Phase-Separation\\
Characterisation for Erd\H{o}s Problem 249\\
{\large A companion note on transport, curvature, and local obstructions}}
\author{Will Cook}
\date{July 2026}

\begin{document}
\sloppy
\maketitle
\begin{abstract}
Let $S=\sum_{n\geq1}\ph(n)/2^n$ and
$R_N=\sum_{j\geq1}\ph(N+j)/2^j$.  The M\"obius residue kernel governing
transport of the tails $R_N$ is affine in its scale coordinate, with periodic
slope and intercept.  Consequently, every finite coefficient system satisfying
$\sum_i c_i=\sum_i c_im_i=0$ annihilates all divisor channels already present
at a base height.  Curvature, prime-jump, balanced $(1,2,3,6)$, and anchored
$(1,3,5,15)$ observables admit exact dyadic-window decompositions and
decidable central-residue conditions implying non-integrality.  A
first-harmonic saving on a finite block forces a tail-difference certificate,
and an exact four-term large-prime pivot decomposition reduces such a saving
to centred, fibre-mean, bad-supplier, and non-supplier budgets.  The fibre-free
cofinal window-pair separation predicate is equivalent to the irrationality
of $S$.  We also prove two local obstructions: every fixed-precision
valuation--unit word has a prefix-locked centred completion, and a synthetic
nonzero lcm carry satisfying the stated factor and anchor constraints survives
every finite integer shift polynomial.  The finite criteria and obstruction
theorems are unconditional.  Erd\H{o}s Problem~249 remains open because no
cofinal arithmetic input or harmonic bound is proved.
\end{abstract}

\section{Main results and logical boundary}\label{sec:scope}

The gateway exposition~\cite{mainpaper} proves the denominator exclusion and
the complete finite-certificate normal form for $S$.  This paper has a
different thesis.  Its argument has three stages.  First, finite differences
annihilate the old affine divisor channels and turn the surviving channels
into exact dyadic-window residues.  Second, a first-harmonic gap turns residue
dispersion into a finite certificate, while the large-prime pivot
decomposition identifies four precise budgets that would produce that gap.
Third, phase separation gives a complete geometric normal form and two
countermodels delimit what fixed local transport data can prove.  Thus the
paper isolates exact implications and exact obstructions; it does not prove
the cofinal arithmetic input needed to settle the problem.

\begin{theorem}[affine-channel annihilation]\label{thm:affine}
Let \(m_i\in\mathbb N\) and \(c_i\in\mathbb R\) satisfy
\(\sum_i c_i=0\) and \(\sum_i c_im_i=0\).  If \(d\mid H\), then
\[
  \sum_i c_i K_d(m_iH)=0,
\]
where \(K_d\) is the exact M\"obius transport-residue kernel.
\end{theorem}
\noindent Formalised as
\lword{JointExponentTransport.lean}{36}{oldChannel_affine_moment_annihilation}{affine-channel annihilation}.
The anchored coefficients at \(H,3H,5H,15H\) are one primitive instance.
Sections~\ref{sec:windows}--\ref{sec:primejump} turn this cancellation into
finite central-residue criteria.

\paragraph{The harmonic bridge.}
Theorem~\ref{thm:first-pivot} decomposes the first-harmonic block sum into
centred, fibre-mean, bad-supplier, and non-supplier terms.  Its exact budget
implies a finite tail-difference certificate; the corresponding cofinal
budget would imply irrationality.  This is the analytic implication connecting
the transport identities to the phase geometry, not an additional proof of
the missing estimates.

\begin{theorem}[phase-separation characterisation]\label{thm:phase}
The cofinal fibre-free window-pair separation predicate of
Section~\ref{sec:pivot} holds if and only if \(S\) is irrational.
\end{theorem}
\noindent Formalised as the
\lword{PivotAntiReconstruction.lean}{1744}{dtwWindowSeparatedPairs_iff_irrational_totient_series}{phase-separation equivalence}.
This is an exact equivalence, not a weaker sufficient condition.  Its required
arithmetic input remains open.

\begin{proposition}[local no-go boundaries]\label{prop:nogo}
At every fixed valuation--unit precision, each compatible finite word has a
prefix-locked centred carry completion.  Separately, for every lcm height
\(t\ge3\), a nonzero synthetic carry satisfies the stated factor ideals,
whole-ray anchors, and every finite integer shift polynomial while remaining
a strict survivor.
\end{proposition}
\noindent Formalised by the
\lword{TropicalCurvatureCarry.lean}{137}{fixedPrecisionTropicalNoGo}{fixed-precision countermodel}
and the
\lword{LcmFactorIdealPulseObstruction.lean}{798}{lcm_factorIdeal_finiteRank_shiftAlgebra_not_sufficient}{finite shift-algebra countermodel}.
These are scoped countermodels, not global impossibility theorems.

Every positive endpoint below still has the form
\[
 \text{cofinal supply of finite certificates}\Longrightarrow S\notin\mathbb Q.
\]
The finite identities, criteria, equivalence, and countermodels are
unconditional; no cofinal supply is proved.  The open problem is recorded in
the original collection and maintained catalogue
\cite{erdosgraham1980,erdos1988,erdosproblems}.  The Lambert-series and
M\"obius/divisor background is classical~\cite{erdos1948,apostol};
Borwein's theorem gives a neighbouring general irrationality
benchmark~\cite{borwein1992}.  None supplies the cofinal input isolated
here.  The theorem-specific comparison boundary is maintained with the
gateway paper.  Linked mathematical phrases open the pinned Lean proofs; the
claim registry and release gate that validate these links are described in
the suite's systems companion~\cite{syspaper}.

\paragraph{Reading map.}
Sections~\ref{sec:windows}--\ref{sec:primejump} build the exact transport and
window algebra.  Sections~\ref{sec:harmonic}--\ref{sec:pivot} give the two
global routes: a sufficient harmonic-budget route and an equivalent
phase-separation route.  Sections~\ref{sec:nogo} and~\ref{sec:factorideal}
then show why bounded local transport data cannot by itself produce either
cofinal hypothesis.

\section{Dyadic windows}\label{sec:windows}
For a finite integer combination $X(H)=\sum_i c_iR_{m_iH}$, truncate each
tail at depth $L$.  Multiplication by $2^L$ gives
\[
 2^LX(H)=A(H,L)+E(H,L),
\]
where $A(H,L)\in\Z$ is computable.  If $|E(H,L)|\le B(H,L)$, then
$X(H)\notin\Z$ whenever
\[
 B(H,L)<A(H,L)\bmod2^L<2^L-B(H,L).
\]
This reusable test is the formal
\lword{PrimeJumpWindow.lean}{53}{notMem_int_of_window_remainder_bound}{central-residue criterion}.
Lean checks the exact decomposition, radius, and concrete modular arithmetic;
an external analytic or arithmetic argument must provide unbounded instances.

\section{Curvature on the lcm cone}\label{sec:curvature}
Set
\[
 Q_H=R_{4H}-3R_{2H}+2R_H.
\]
Its depth-$L$ window $C(H,L)$ satisfies
\[
 2^LQ_H=C(H,L)+E_C(H,L),\qquad |E_C(H,L)|\le6H+3L+3.
\]
The exact split and the tail bound are both checked in Lean:
\lword{CurvatureCarry.lean}{119}{two_pow_mul_curvature_eq_window_add_shifted}{exact curvature split}
and
\lword{CurvatureCarry.lean}{137}{abs_curvatureShiftedTail_le}{curvature tail bound}.
The
\lword{CurvatureCarry.lean}{159}{curvature_notMem_int_of_sharpCurvatureCert}{sharp curvature criterion}
therefore proves $Q_H\notin\Z$.  An unbounded supply of these certificates
would imply irrationality by the
\lword{CurvatureCarry.lean}{206}{irrational_totientSeries_of_sharpCurvatureSupply}{curvature-supply implication},
but that supply is not proved.

The window evolves by
$C(H,L+1)=2C(H,L)+d(H,L+1)$; see the
\lword{CurvatureCarry.lean}{259}{curvatureWindow_succ}{exact curvature recurrence}.
Exact phase-certificate and common-carrier formulas expose possible routes to
a supply theorem without asserting one.

\section{Exponent transport and affine cancellation}\label{sec:transport}
If $p\mid H$, multiplication by $p$ adds no squarefree divisor support:
\[
 d\mid pH\Longleftrightarrow d\mid H\qquad(d\text{ squarefree}).
\]
This is the formal
\lword{ExponentOnlyTransport.lean}{36}{squarefree_dvd_mul_iff_of_dvd}{squarefree support invariance}.
The exact residue kernel is affine, not merely periodic:
\[
 K_d(N)=N\,a_d(N)+b_d(N),
\]
where both $a_d$ and $b_d$ are periodic modulo $d$, as recorded in the
\lword{ExponentOnlyTransport.lean}{141}{transportResidueKernel_eq_affineState}{affine transport-kernel theorem}.
Keeping both blocks prevents an intercept-only matrix calculation from being
misread as global non-cancellation. The exact $H\mapsto3H$ two-block transfer
is also formalised
\lword{ExponentOnlyTransport.lean}{176}{residueTransportDefect_three_eq_affineState}{by the three-scale transfer identity}.

Every old channel $d\mid H$ is killed by the first-order defect
\lword{ExponentOnlyTransport.lean}{112}{residueTransportDefect_eq_zero_of_dvd}{in the old-channel cancellation theorem}.
The resulting irrationality theorem remains conditional on a fixed-three
tower supply; see the
\lword{ExponentOnlyTransport.lean}{205}{irrational_totient_series_of_exponentOnlyThreeTransportSupply}{exponent-transport supply implication}.

More generally, coefficients satisfying $\sum_i c_i=\sum_i c_im_i=0$
annihilate every old affine channel
\lword{JointExponentTransport.lean}{36}{oldChannel_affine_moment_annihilation}{by affine moment annihilation}.
For $q(X,Y)=XY-3X-2Y+4$, the vanishing $q(1,1)=q(3,5)=0$ gives
\[
 K_d(15H)-3K_d(3H)-2K_d(5H)+4K_d(H)=0\quad(d\mid H),
\]
as checked in Lean
\lword{JointExponentTransport.lean}{71}{joint35_oldChannel_zero}{by anchored \((3,5)\) cancellation}.
The associated window radius is $19H+5L+5$; it gives a finite
non-integrality criterion
\lword{JointExponentTransport.lean}{172}{joint35Cone_notMem_int_of_cert}{in the joint-transport certificate}.
The corresponding irrationality endpoint is conditional on the
\lword{JointExponentTransport.lean}{217}{irrational_totient_series_of_anchoredJoint35ConeSupply}{anchored joint-cone supply}.

\section{Prime jumps and three-transport boundaries}\label{sec:primejump}
When $p\nmid H$, a jump from $H$ to $pH$ creates new squarefree support.  The
residue Hecke defect splits into a post-jump channel and an explicit migration
correction
\lword{PrimeJumpMigration.lean}{38}{residueHeckeDefect_eq_postJump_add_migration}{by the prime-jump migration identity}.
For $d\mid H$ the total old-channel defect is zero
\lword{PrimeJumpMigration.lean}{46}{residueHeckeDefect_eq_zero_of_dvd}{by prime-jump old-channel cancellation}; the
$(d,H,p)=(5,12,5)$ calculation records the two cancelling terms.

Set
\[
 J(H,p)=(R_{2pH}-R_{pH})-p(R_{2H}-R_H).
\]
Its exact remainder radius is $3pH+(p+1)(L+2)$, with a checked window split
\lword{PrimeJumpWindow.lean}{91}{primeJumpTailCommutator_eq_window_add_remainder}{in the prime-jump window identity}.
This yields a finite non-integrality criterion
\lword{PrimeJumpWindow.lean}{178}{primeJumpTailCommutator_notMem_int_of_sharpKill}{from a sharp prime-jump certificate}.
Lean verifies the instance $(H,p,L)=(12,5,15)$
\lword{PrimeJumpWindow.lean}{186}{primeJumpSharpKill_twelve_five}{in the \((12,5,15)\) calculation}; this is one
verified finite instance, not an unbounded sequence.

The balanced vertices $H,2H,3H,6H$ give a related radius $9H+4L+4$.
The corresponding finite criterion is formalised
\lword{ThreeTransportBoundary.lean}{126}{threeTransport_notMem_int_of_cert}{as the three-transport certificate};
the irrationality theorem is conditional on an unbounded supply
\lword{ThreeTransportBoundary.lean}{177}{irrational_totientSeries_of_roughnessBoundaryLatticeSupply}{through the roughness-boundary implication}.
The concrete instance $(H,L)=(60,12)$ is checked
\lword{ThreeTransportBoundary.lean}{203}{sharpThreeTransportCert_60_12}{in the \((60,12)\) calculation}.

\section{A first-harmonic criterion}\label{sec:harmonic}
For the window discrepancy $A(h,N,L)$ define
\[
 W_1(h,N,L)=\cos\left(2\pi\frac{A(h,N,L)\bmod2^L}{2^L}\right).
\]
Assume $N\in[X,2X)$ and $16(2X+h+L+2)\le2^L$.  If every $N$ failed the
central-residue certificate, each phase would lie within $\pi/8$ of an
integer turn, so each cosine would exceed $9/10$.  Hence
\[
 \sum_{X\le N<2X}W_1(h,N,L)\le\frac9{10}X
\]
forces a certificate in the block
\lword{FirstHarmonicGap.lean}{128}{exists_certifiedKill_of_first_harmonic_gap}{by the block first-harmonic criterion}.
More generally, if $T$ is any nonempty finite set of indices below $2X$ and
\[
 \sum_{N\in T}W_1(h,N,L)\le\frac9{10}|T|,
\]
then a certificate occurs at some $N\in T$
\lword{FirstHarmonicGap.lean}{104}{exists_certifiedKill_of_first_harmonic_gap_subset}{by the finite-subset criterion}.
By certificate completeness, the witness produced here is an exact
non-integral tail difference, not a heuristic phase test.  Hence a
first-harmonic saving on one finite block produces exactly the kind of witness
required by the tail-difference criterion.  What is still missing is a
first-harmonic estimate that supplies such blocks at unbounded parameters; no
such estimate is proved here.

\section{The exact large-prime pivot decomposition}\label{sec:first-pivot}
The preceding criterion becomes useful only after the harmonic sum has been
split into terms for which an arithmetic estimate is plausible.  We record
that split explicitly.  Put
\[
 Z_N=\exp\!\left(
   2\pi i\,\frac{A(h,N,L)\bmod 2^L}{2^L}\right),
 \qquad X\le N<2X,
\]
so that \(\Re Z_N=W_1(h,N,L)\) and \(|Z_N|=1\).
Fix an overlap depth \(s\), set \(q_N=N+(L-s+1)\), and write
\(q_N=m_Np_N\), where \(p_N\) is its largest prime factor.  Call \(N\) a
\emph{supplier} when
\[
 0<m_N\le\frac{\sqrt X}{2},
 \qquad 2\sqrt X<p_N,
\]
with the displayed factorisation certified; call it \(\eta\)-good when also
\(\eta m_N\le\ph(m_N)\).  On a supplier define the pivot phase
\[
 P_N=\exp\!\left(
 2\pi i\,\frac{(2^h-1)\ph(m_N)(p_N-1)}{2^{L-s+1}}\right),
 \qquad \rho_N=\frac{Z_N}{P_N}.
\]
The identity \(\ph(m_Np_N)=\ph(m_N)(p_N-1)\) is exact here because the size
conditions force \(p_N\nmid m_N\).  For the supplier fibre
\(\mathcal F_m=\{N:m_N=m\}\), let
\[
 \overline P_m=\frac1{|\mathcal F_m|}
                 \sum_{N\in\mathcal F_m}P_N.
\]
Finally, let \(\mathcal G,\mathcal B,\mathcal U\) be respectively the good
suppliers, bad suppliers, and non-suppliers in \([X,2X)\), and put
\[
\begin{aligned}
 C&=\sum_{N\in\mathcal G}\rho_N(P_N-\overline P_{m_N}),&
 M&=\sum_{N\in\mathcal G}\rho_N\overline P_{m_N},\\
 B&=\sum_{N\in\mathcal B}Z_N,&
 U&=\sum_{N\in\mathcal U}Z_N.
\end{aligned}
\]

\begin{theorem}[four-term pivot decomposition and budget]
\label{thm:first-pivot}
For every choice of \(h,X,L,s,\eta\),
\[
  \sum_{X\le N<2X}Z_N=C+M+B+U.
\]
If \(X>0\), \(16(2X+h+L+2)\le2^L\), and
\[
 \Re C\le\frac{14}{25}X,\qquad
 |M|\le\frac1{100}X,\qquad
 |B|\le\frac1{100}X,\qquad
 |U|\le\frac8{25}X,
\]
then
\[
 \sum_{X\le N<2X}W_1(h,N,L)\le\frac9{10}X,
\]
so the block contains a finite tail-difference certificate.  Consequently,
if for every \(h>0\) there are fixed \(s>0\) and \(0<\eta<1\) for which these
four budgets hold at arbitrarily large \(X\), with \(h\le L-s\), then \(S\)
is irrational.
\end{theorem}
\noindent The exact split is formalised
\lword{FirstHarmonicPivot.lean}{342}{windowFirstExp_sum_eq_pivot_decomposition}{by the four-term decomposition};
the budget implication is
\lword{FirstHarmonicPivot.lean}{377}{first_harmonic_gap_of_pivotBudgetAt}{the one-sided pivot budget};
and the cofinal conditional endpoint is
\lword{FirstHarmonicPivot.lean}{405}{irrational_totient_series_of_pivotResidualDecorrelation}{the residual-decorrelation implication}.

The constants are not heuristic bookkeeping:
\[
 \frac{14}{25}+\frac1{100}+\frac1{100}+\frac8{25}
 =\frac9{10}.
\]
Only the real part of the centred term is bounded; no unnecessary complex-norm
estimate is assumed there.  The theorem proves the exact algebraic reduction.
It does not prove a prime-distribution or residual-decorrelation estimate
supplying the four inequalities at cofinal scales.

\section{Adjacent-pair phase separation}\label{sec:adjacentphase}
At a scale satisfying the same room inequality, with $N+1<2X$, suppose that
the first harmonic phases at two adjacent indices obey
\[
 \left\lVert
 \operatorname{windowFirstExp}(h,N+1,L)-\operatorname{windowFirstExp}(h,N,L)
 \right\rVert^2\ge\frac{19}{25}.
\]
Then a certificate occurs at one of $N$ and $N+1$
\lword{AdjacentPhaseSeparation.lean}{258}{exists_certifiedKill_of_phasePairSeparation}{by adjacent-phase separation}.
Consequently, a cofinal supply of adjacent pairs satisfying this inequality
implies irrationality of the totient constant
\lword{AdjacentPhaseSeparation.lean}{295}{irrational_totient_series_of_adjacentPhaseSeparation}{through the adjacent-phase supply implication}.
No cofinal adjacent-pair separation is proved here.

\section{Pivot anti-reconstruction and window-pair geometry}\label{sec:pivot}
The first-harmonic phase of a finite window is the exact additive character of
its dyadic discrepancy residue.  It differs from the corresponding infinite
tail-orbit phase by the explicitly bounded omitted tail.  On a finite supplier
fibre, the anchored defect is the squared distance of the window phases from
the constant phase $1$; its centred form is the usual pairwise-energy identity.
Thus a counted family of separated phase pairs of sufficient total squared
chord mass produces a finite certificate.

The formal proof records several equivalent sufficient cofinal hypotheses.  In
particular, a cofinal pivot
anti-reconstruction bound, a cofinal centred reconstruction-variance bound, or
a cofinal counted-separated-pairs bound is sufficient for irrationality of
$S$. The direct finite criterion and the resulting conditional irrationality
theorem are both formalised
\lword{PivotAntiReconstruction.lean}{1105}{exists_certifiedKill_of_pivotFiberSeparatedPairs}{by the pivot-fibre certificate}
and
\lword{PivotAntiReconstruction.lean}{1502}{irrational_totient_series_of_pivotAntiReconstruction}{the pivot anti-reconstruction implication}.

There is also an exact fibre-free formulation: the cofinal window-pair
predicate is equivalent to irrationality of $S$, not merely sufficient for it.
The reverse implication uses the expansive doubling orbit to obtain adjacent
infinite-tail separation and then chooses a finite window depth
\lword{PivotAntiReconstruction.lean}{1744}{dtwWindowSeparatedPairs_iff_irrational_totient_series}{in the exact window-pair equivalence}.
The development does not prove any of these cofinal arithmetic hypotheses, so it does
not settle Erd\H{o}s~\#249.

\section{Fixed local precision cannot finish the argument}\label{sec:nogo}
A finite valuation--unit symbol records a bounded amount of local transport
data.  For every positive precision $u$, every finite compatible word admits
a prefix-locked centred completion.  Thus no finite word over this fixed local
alphabet excludes all centred endpoints
\lword{TropicalCurvatureCarry.lean}{137}{fixedPrecisionTropicalNoGo}{by the fixed-precision no-go theorem}.
This does not say curvature certificates are rare.  A proof may use precision
growing with scale, global correlations, or a harmonic gap.  It rules out only
the inference that a bounded local signature alone forces escape.

\section{LCM factor ideals and whole-ray anchors}\label{sec:factorideal}
For every $t\geq 3$, there is a nonzero synthetic affine carry whose forcing
letters lie in the principal ideal generated by
$\varphi(\operatorname{lcm}(1,\ldots,t))$, satisfy the exact whole-ray anchor
values, and remain inside the relevant diagonal bounds. Nevertheless its
state is a strict survivor
\lword{LcmFactorIdealPulseObstruction.lean}{866}{lcm_factorIdeal_sparseAnchor_not_sufficient}{in the sparse-anchor countermodel}.

The same example remains a dyadic coboundary after every finite integer shift
polynomial. In particular, the finite-rank LCM shift algebra alone cannot
remove its terminal compensation
\lword{LcmFactorIdealPulseObstruction.lean}{798}{lcm_factorIdeal_finiteRank_shiftAlgebra_not_sufficient}{in the finite shift-algebra countermodel}.
This is a synthetic countermodel: it does not assert that these forcing letters
are actual totient differences, and it does not rule out a nonlinear argument
using the missing fresh divisor channels.

\section{Conclusion and exact boundary}\label{sec:index}

The results form one chain rather than a list of unrelated certificate
constructions.  Affine transport explains which divisor channels a finite
difference removes; a dyadic window converts the remaining tail observable
into an exact central-residue test; and harmonic or pairwise phase dispersion
turns that test into a certificate.  The four-term pivot theorem makes the
first of these global routes arithmetically explicit, while the
window-pair theorem shows that the second is logically complete.

What remains missing is not another finite identity.  One must produce the
required dispersion cofinally, either through estimates meeting the four
pivot budgets or through an equivalent window-pair separation mechanism.
The two countermodels explain why fixed local precision, factor ideals,
whole-ray anchors, and finite shift algebra do not force such dispersion.
They delimit the current methods without asserting a global impossibility.

For reference, the exact status of each layer is:
\begin{center}\small
\begin{tabular}{@{}
>{\raggedright\arraybackslash}p{0.21\textwidth}
>{\raggedright\arraybackslash}p{0.40\textwidth}
>{\raggedright\arraybackslash}p{0.29\textwidth}@{}}
\toprule Layer & Checked statement & Status \\
\midrule
Curvature & exact split, radius, recurrence, criterion & finite certificate criterion \\
Exponent transport & stable support, affine state, cancellation & exact identity \\
Joint $(3,5)$ transport & moment annihilation and criterion & identity and finite criterion \\
Prime jump & migration, sharp window, $(12,5,15)$ & identity and finite instance \\
Three transport & sharp window and $(60,12)$ & criterion and finite instance \\
First harmonic & cosine gap forces a certificate & finite certificate criterion \\
Large-prime pivot & exact four-term split and one-sided budget & identity and conditional endpoint \\
Adjacent phases & squared chord separation forces a certificate & finite certificate criterion \\
Pivot anti-reconstruction & pairwise phase energy and exact window-pair equivalence & exact equivalence \\
Tropical signature & fixed precision has a completion & scoped countermodel \\
LCM factor ideals & factor/anchor and finite shift algebra have a synthetic survivor & scoped synthetic countermodel \\
Irrationality of $S$ & follows from a named unbounded supply & open \\
\bottomrule
\end{tabular}
\end{center}
No cofinal arithmetic input or harmonic bound is proved.  Therefore the chain
stops at exact finite criteria and equivalent reformulations, and
Erd\H{o}s~\#249 remains open.

\begin{thebibliography}{9}
\bibitem{erdos1948} P.~Erd\H{o}s,
\emph{On arithmetical properties of Lambert series},
J. Indian Math. Soc. (N.S.) \textbf{12} (1948), 63--66,
\href{https://users.renyi.hu/~p_erdos/1948-04.pdf}{archive scan}.
\bibitem{apostol} T.~M. Apostol,
\emph{Introduction to Analytic Number Theory},
Undergraduate Texts in Mathematics, Springer, 1976,
doi:\href{https://doi.org/10.1007/978-1-4757-5579-4}{10.1007/978-1-4757-5579-4}.
\bibitem{borwein1992} P.~B. Borwein,
\emph{On the irrationality of certain series},
Math. Proc. Cambridge Philos. Soc. \textbf{112} (1992), no.~1, 141--146,
doi:\href{https://doi.org/10.1017/S030500410007081X}{10.1017/S030500410007081X}.
\bibitem{erdosgraham1980} P.~Erd\H{o}s and R.~L. Graham,
\emph{Old and New Problems and Results in Combinatorial Number Theory},
Monographies de l'Enseignement Math\'ematique 28 (1980), 61--62,
\href{https://mathweb.ucsd.edu/~ronspubs/80_11_number_theory.pdf}{scan}.
\bibitem{erdos1988} P.~Erd\H{o}s,
\emph{On the irrationality of certain series: problems and results},
in \emph{New Advances in Transcendence Theory}, ed.~A.~Baker,
Cambridge University Press, 1988, pp.~102--109,
doi:\href{https://doi.org/10.1017/CBO9780511897184.009}{10.1017/CBO9780511897184.009}.
\bibitem{erdosproblems} T.~Bloom (ed.), \emph{Erd\H{o}s Problems}, entry 249,
\href{https://www.erdosproblems.com/249}{\texttt{erdosproblems.com/249}},
accessed July 2026.
\bibitem{mainpaper} W.~Cook, \emph{Tail Certificates and Achievement-Set
Geometry for Erd\H{o}s Problems 249 and 257}, gateway exposition in this
release.
\bibitem{syspaper} W.~Cook, \emph{Claim-Faithful Publication of a Formalised
Mathematics Corpus}, systems companion in this release.
\end{thebibliography}
\end{document}
